\begin{lemma}[The differential structure conditions hold at every point]
Let $n\ge2$ be an integer, let $p>1$, $\varepsilon\ge0$, $\gamma_0>0$ and
$\gamma_1>0$ be real, let $\Omega\subseteq\R^n$ be open, let
$A\colon\R^n\times\R\times\R^n\to\R^n$ and let $\sigma\colon\R\to\R$.  Assume
\begin{itemize}
\item $A$ satisfies (Q1) of the thesis--(Q2) of the thesis on $\Omega$ with the parameters
  $p,\varepsilon,\gamma_0,\gamma_1$ (\noderef{EllipticStructureQ12}); and
\item $A$ satisfies (Q6) of the thesis on $\Omega$ with modulus $\sigma$
  (\noderef{EllipticStructureQ6}).
\end{itemize}
Then the two differential inequalities of (Q1) of the thesis--(Q2) of the thesis hold at
\emph{every} point of $\Omega$: for every $x\in\Omega$, every $z\in\R$ and every
$h\in\R^n$ with $\varepsilon+|h|^2\ne0$,
\[
\gamma_0\bigl(\varepsilon+|h|^2\bigr)^{\frac p2-1}|\xi|^2
\ \le\ \sum_{i,j=1}^nD_{h_j}A^i(x,z,h)\,\xi^i\xi^j
\qquad\text{for every }\xi\in\R^n ,
\tag{a}
\]
\[
\bigl|D_{h_j}A^i(x,z,h)\bigr|
\ \le\ \gamma_1\bigl(\varepsilon+|h|^2\bigr)^{\frac p2-1}
\qquad\text{for all }i,j\in\{1,\dots,n\} .
\tag{b}
\]
Here $D_{h_j}A^i(x,z,h)$ is the $i$-th component of the derivative of
$h\mapsto A(x,z,h)$ at $h$ in the direction of the $j$-th standard basis vector,
and the sum in (a) is the value of that derivative at $\xi$ paired with $\xi$,
exactly as in \noderef{EllipticStructureQ12}.

\medskip
\noindent\textbf{What this lemma is, and what it is not.}  The paper states
(Q1) of the thesis and (Q2) of the thesis ``for all $h,\xi\in\R^n$, all $z\in\R$ and
a.e.\ $x\in\Omega$'' (p.~1), and \noderef{EllipticStructureQ12} transcribes that
almost-everywhere quantifier faithfully.  This lemma does \emph{not} strengthen
that hypothesis: it \emph{derives} the pointwise form from it \emph{together
with} (Q6) of the thesis, which is an additional hypothesis, stated pointwise in $x$
by the paper, and which \S5 assumes throughout.  So the content of the lemma is
the implication
\[
\text{(Q1) of the thesis--(Q2) of the thesis a.e.}\ +\ \text{(Q6) of the thesis pointwise}
\ \Longrightarrow\ \text{(Q1) of the thesis--(Q2) of the thesis pointwise},
\]
and nothing is claimed about an operator for which (Q6) of the thesis fails.

\noindent\textbf{Why it is needed.}  Section 5 freezes the operator at a point:
for $x_0\in\Omega'$ and $R\le R_0$ it solves the Dirichlet problem
(5.2) of the thesis for the field $h\mapsto A(x_0,u(x_0),h)$ (p.~46).  That field is
the restriction of $A$ to the \emph{single} point $x_0$, so every result applied
to it --- Lemma~Lemma~3.1 of the thesis and its cited input (C3), and Theorem~1 ---
requires (Q1) of the thesis--(Q2) of the thesis \emph{at} $x_0$, and the centres $x_0$
range over all of a subdomain, not over a set of full measure: the comparison
scheme behind Theorem~5 consumes the comparison data at every centre of an open
set.  An almost-everywhere hypothesis therefore does not suffice, and this
lemma is what closes the gap.

\noindent\textbf{The constants are not degraded.}  The conclusion carries the
same $\gamma_0$ and $\gamma_1$ as the hypothesis, not $\gamma_0$ and $\gamma_1$
spoiled by a factor involving $\sigma$.  This matters because the consumers
instantiate the cited results at the paper's constants: the frozen operator of
\S5 must satisfy (Q1) of the thesis--(Q2) of the thesis with $\gamma_0,\gamma_1$ for
Theorem~1's constant to be the paper's $c_0=c_0(p,\gamma_1/\gamma_0,n)$.  The
reason no loss occurs is that $\sigma$ enters only through a limit in which it
tends to $0$: it is used to transport an inequality from nearby points, never to
weaken it.
\end{lemma}
\begin{proof}
Write
\[
w(y)=\varepsilon+|y|^2\quad(y\in\R^n),
\qquad
D=\{y\in\R^n:w(y)\ne0\} .
\]
Since $\varepsilon\ge0$, $w\ge0$ everywhere, so $w>0$ on $D$; and $w$ is
continuous, so $D$, the preimage of the open set $\R\setminus\{0\}$, is open.

\emph{Step 0: the good points, and the pointwise differentiability.}  The clause
(Q1) of the thesis of \noderef{EllipticStructureQ12} is an almost-everywhere
statement: there is a Lebesgue null set $N_1$ such that every $x\notin N_1$ with
$x\in\Omega$ satisfies
\[
\gamma_0w(h)^{\frac p2-1}|\xi|^2\le\sum_{i,j}D_{h_j}A^i(x,z,h)\xi^i\xi^j
\qquad\text{for all }z\in\R,\ h\in D,\ \xi\in\R^n .
\tag{Q1$_x$}
\]
Likewise the clause (Q2) of the thesis gives a null set $N_2$ such that every
$x\notin N_2$ with $x\in\Omega$ satisfies
\[
\bigl|D_{h_j}A^i(x,z,h)\bigr|\le\gamma_1w(h)^{\frac p2-1}
\qquad\text{for all }z\in\R,\ h\in D,\ i,j .
\tag{Q2$_x$}
\]
Put $N=N_1\cup N_2$, a Lebesgue null set, and
\[
G=\Omega\setminus N\qquad(\text{the \emph{good} points}),
\]
so every $x\in G$ satisfies both (Q1$_x$) and (Q2$_x$).

We also record what holds at \emph{every} point.  By the clause
\texttt{hdiff} of \noderef{EllipticStructureQ12}, for every $x\in\Omega$ and
every $z\in\R$ the map $A(x,z,\cdot)$ is of class $C^1$ on $D$.  As $D$ is open,
this gives for each $h\in D$ that $A(x,z,\cdot)$ is (Fr\'echet) differentiable at
$h$, with derivative the linear map $\xi\mapsto D_\xi A(x,z,h)$ whose $i$-th
component at $\xi=e_j$ is the $D_{h_j}A^i(x,z,h)$ of the statement, and that
$h\mapsto D_\cdot A(x,z,h)$ is continuous on $D$.

\emph{Step 1: every point of $\Omega$ is a limit of good points.}  Let
$x_1\in\Omega$.  Since $\Omega$ is open there is $\delta_0>0$ with
$B_{\delta_0}(x_1)\subseteq\Omega$.  For $k\in\mathbb{N}$ put
$\delta_k=\min\{\delta_0,\tfrac1{k+1}\}>0$.  The ball $B_{\delta_k}(x_1)$ has
measure $\omega_n\delta_k^{\,n}>0$, where $\omega_n=|B_1(0)|\in(0,\infty)$ (the
Preamble item on the volume of a ball), while $N$ is null; so
$B_{\delta_k}(x_1)\not\subseteq N$ and we may choose
\[
x_k\in B_{\delta_k}(x_1)\setminus N\subseteq\Omega\setminus N=G .
\]
Then $x_k\in G$ for every $k$ and $|x_k-x_1|<\delta_k\le\frac1{k+1}\to0$, i.e.
$x_k\to x_1$.

\emph{Step 2: admissible segments and the weight along them.}  Let $h\in D$ and
$\eta\in\R^n$.  The map $s\mapsto h+s\eta$ is continuous and sends $0$ to
$h\in D$ with $D$ open, so there is $t_0>0$ --- depending only on $h,\eta$ and
$\varepsilon$, and in particular not on $x$ --- with
\[
h+s\eta\in D\qquad\text{for all }|s|\le t_0 .
\]
On $[0,t_0]$ the function $s\mapsto w(h+s\eta)$ is continuous with values in
$(0,\infty)$, and $u\mapsto u^{\frac p2-1}$ is continuous on $(0,\infty)$, so
\[
W_\eta(s)=w(h+s\eta)^{\frac p2-1}
\]
is continuous and positive on $[0,t_0]$; being continuous on a compact interval
it is interval-integrable there.

\emph{Step 3: the integrated form of (Q1) at a good point.}  Let $x\in G$,
$z\in\R$, $h\in D$ and $\xi\in\R^n$, let $t_0$ be as in Step 2 for
$\eta=\xi$, and let $0<t\le t_0$.  Define
\[
\varphi(s)=\sum_{i=1}^nA^i(x,z,h+s\xi)\,\xi^i\qquad(|s|\le t_0).
\]
For such $s$ the point $h+s\xi$ lies in $D$, where $A(x,z,\cdot)$ is
differentiable (Step 0); composing with the affine map $s\mapsto h+s\xi$, whose
derivative is the constant vector $\xi$, and then with the fixed linear
functional $v\mapsto\sum_iv^i\xi^i$, we get that $\varphi$ is differentiable at
$s$ with
\[
\varphi'(s)=\sum_{i,j}D_{h_j}A^i(x,z,h+s\xi)\,\xi^i\xi^j .
\]
By Step 0 the derivative map is continuous on $D$, so $\varphi'$ is continuous
on $[0,t]$ and hence interval-integrable there.  The fundamental theorem of
calculus in interval-integral form (Preamble) therefore gives
\[
\varphi(t)-\varphi(0)=\int_0^t\varphi'(s)\,ds .
\]
By (Q1$_x$), applied at the point $h+s\xi\in D$ and the direction $\xi$,
\[
\varphi'(s)\ \ge\ \gamma_0|\xi|^2W_\xi(s)\qquad(0\le s\le t),
\]
and both sides are interval-integrable on $[0,t]$ (Step 2 for the right-hand
side), so monotonicity of the interval integral (Preamble) yields
\[
\sum_{i=1}^n\bigl(A^i(x,z,h+t\xi)-A^i(x,z,h)\bigr)\xi^i
\ \ge\ \gamma_0|\xi|^2\int_0^tW_\xi(s)\,ds
\qquad(x\in G,\ 0<t\le t_0).
\tag{1}
\]

\emph{Step 4: the integrated form of (Q2) at a good point.}  Let $x\in G$,
$z\in\R$, $h\in D$ and $i,j\in\{1,\dots,n\}$, let $e_j$ be the $j$-th standard
basis vector, let $t_0$ be as in Step 2 for $\eta=e_j$ and let $0<t\le t_0$.
Define $\psi(s)=A^i(x,z,h+se_j)$ for $|s|\le t_0$.  Exactly as in Step 3,
$\psi$ is differentiable with
$\psi'(s)=D_{h_j}A^i(x,z,h+se_j)$, and $\psi'$ is continuous, hence
interval-integrable, on $[0,t]$.  By (Q2$_x$) at the point $h+se_j\in D$,
\[
|\psi'(s)|\le\gamma_1W_{e_j}(s)\qquad(0\le s\le t).
\]
The fundamental theorem of calculus (Preamble) gives
$\psi(t)-\psi(0)=\int_0^t\psi'(s)\,ds$, so by the norm inequality for the
integral and then monotonicity (Preamble, parts (c) and (b)),
\[
\bigl|A^i(x,z,h+te_j)-A^i(x,z,h)\bigr|
\ \le\ \int_0^t|\psi'(s)|\,ds
\ \le\ \gamma_1\int_0^tW_{e_j}(s)\,ds
\qquad(x\in G,\ 0<t\le t_0).
\tag{2}
\]

\emph{Step 5: transport to an arbitrary point by (Q6) of the thesis.}  Fix
$x_1\in\Omega$ and let $(x_k)\subseteq G$ with $x_k\to x_1$ be the sequence of
Step 1.  The clause (Q6) of the thesis of \noderef{EllipticStructureQ6} holds for
every pair of points of $\Omega$, in particular for $(x_k,x_1)$, and with
$z'=z$ it reads
\[
\bigl|A(x_k,z,y)-A(x_1,z,y)\bigr|
\ \le\ \sigma\bigl(|x_k-x_1|\bigr)\bigl(1+|y|^2\bigr)^{\frac{p-1}2}
\qquad(y\in\R^n).
\]
The clauses $\sigma$-continuity and $\sigma(0)=0$ of the same node, together
with $|x_k-x_1|\to0$, give $\sigma(|x_k-x_1|)\to\sigma(0)=0$; hence, for each
fixed $y\in\R^n$, the factor $(1+|y|^2)^{(p-1)/2}$ being a finite constant,
\[
A(x_k,z,y)\longrightarrow A(x_1,z,y)\qquad(k\to\infty).
\tag{3}
\]
Now fix $z$, $h\in D$, $\xi$ and $0<t\le t_0$ as in Step 3.  Applying $(3)$ at
the two points $y=h+t\xi$ and $y=h$ and taking the (finite) linear combination
with the fixed coefficients $\xi^i$, the left-hand side of $(1)$ at $x=x_k$
converges to the same expression at $x_1$, while the right-hand side of $(1)$
does not depend on $x$ at all.  A non-strict inequality between convergent
sequences of reals passes to the limit, so
\[
\sum_{i=1}^n\bigl(A^i(x_1,z,h+t\xi)-A^i(x_1,z,h)\bigr)\xi^i
\ \ge\ \gamma_0|\xi|^2\int_0^tW_\xi(s)\,ds
\qquad(0<t\le t_0).
\tag{1$'$}
\]
The same argument applied to $(2)$ --- using $(3)$ at $y=h+te_j$ and $y=h$,
the componentwise bound $|v^i|\le|v|$ to pass from the convergence of the
vectors to that of their $i$-th components, and the continuity of the absolute
value --- gives
\[
\bigl|A^i(x_1,z,h+te_j)-A^i(x_1,z,h)\bigr|
\ \le\ \gamma_1\int_0^tW_{e_j}(s)\,ds
\qquad(0<t\le t_0).
\tag{2$'$}
\]
Note that the sequence here is a sequence of \emph{points} $x_k$ of $\Omega$ and
that $t$, $h$, $\xi$, $i$, $j$ are held fixed while $k\to\infty$; no limit of
functions and no almost-everywhere convergence is involved.

\emph{Step 6: the average of the weight.}  Let $h\in D$, $\eta\in\R^n$ and
$t_0$ as in Step 2.  Then
\[
\frac1t\int_0^tW_\eta(s)\,ds\ \longrightarrow\ W_\eta(0)=w(h)^{\frac p2-1}
\qquad(t\downarrow0).
\tag{4}
\]
Indeed, let $\theta>0$.  By continuity of $W_\eta$ at $0$ (Step 2) there is
$\delta\in(0,t_0]$ with $|W_\eta(s)-W_\eta(0)|\le\theta$ for $0\le s\le\delta$.
For $0<t\le\delta$, using $\int_0^tW_\eta(0)\,ds=tW_\eta(0)$, the norm
inequality for the integral and monotonicity (Preamble),
\[
\Bigl|\frac1t\int_0^tW_\eta(s)\,ds-W_\eta(0)\Bigr|
=\frac1t\Bigl|\int_0^t\bigl(W_\eta(s)-W_\eta(0)\bigr)ds\Bigr|
\le\frac1t\int_0^t\bigl|W_\eta(s)-W_\eta(0)\bigr|ds\le\theta .
\]

\emph{Step 7: differentiating back, and the two conclusions.}  Fix $x_1\in\Omega$,
$z\in\R$ and $h\in D$; recall from Step 0 that $A(x_1,z,\cdot)$ is
differentiable at $h$.

\emph{(a)}  Let $\xi\in\R^n$.  If $\xi=0$ then both sides of (a) vanish --- the
derivative is a linear map, so its value at $0$ is $0$, and $|\xi|^2=0$ --- and
(a) holds.  So let $\xi\ne0$, let $t_0$ be as in Step 2 for $\eta=\xi$ and
divide $(1')$ by $t>0$:
\[
\frac1t\sum_{i=1}^n\bigl(A^i(x_1,z,h+t\xi)-A^i(x_1,z,h)\bigr)\xi^i
\ \ge\ \gamma_0|\xi|^2\cdot\frac1t\int_0^tW_\xi(s)\,ds
\qquad(0<t\le t_0).
\]
As $t\downarrow0$ the left-hand side converges to
$\sum_{i,j}D_{h_j}A^i(x_1,z,h)\xi^i\xi^j$: differentiability of
$A(x_1,z,\cdot)$ at $h$ means precisely that
$\bigl(A(x_1,z,h+t\xi)-A(x_1,z,h)\bigr)/t\to D_\xi A(x_1,z,h)$ in $\R^n$ as
$t\to0$, and pairing with the fixed vector $\xi$ is a continuous operation,
while $\sum_iD_\xi A^i(x_1,z,h)\xi^i=\sum_{i,j}D_{h_j}A^i(x_1,z,h)\xi^i\xi^j$
by linearity of the derivative in the direction.  The right-hand side converges
to $\gamma_0|\xi|^2w(h)^{\frac p2-1}$ by $(4)$.  A non-strict inequality between
two convergent families passes to the limit, so
\[
\gamma_0w(h)^{\frac p2-1}|\xi|^2\le\sum_{i,j}D_{h_j}A^i(x_1,z,h)\xi^i\xi^j ,
\]
which is (a) at $x_1$, with the constant $\gamma_0$ of the hypothesis.

\emph{(b)}  Let $i,j\in\{1,\dots,n\}$, let $t_0$ be as in Step 2 for
$\eta=e_j$ and divide $(2')$ by $t>0$:
\[
\frac1t\bigl|A^i(x_1,z,h+te_j)-A^i(x_1,z,h)\bigr|
\ \le\ \gamma_1\cdot\frac1t\int_0^tW_{e_j}(s)\,ds
\qquad(0<t\le t_0).
\]
As $t\downarrow0$ the left-hand side converges to
$|D_{h_j}A^i(x_1,z,h)|$, by the same difference-quotient convergence as in (a)
--- now in the direction $e_j$ --- followed by the continuity of
$v\mapsto|v^i|$ on $\R^n$; and the right-hand side converges to
$\gamma_1w(h)^{\frac p2-1}$ by $(4)$.  Passing to the limit gives
\[
\bigl|D_{h_j}A^i(x_1,z,h)\bigr|\le\gamma_1w(h)^{\frac p2-1},
\]
which is (b) at $x_1$, with the constant $\gamma_1$ of the hypothesis.

Since $x_1\in\Omega$, $z\in\R$, $h\in D$ and the remaining data were arbitrary,
(a) and (b) hold at every point of $\Omega$, and the proof is complete.
\end{proof}
